\documentclass[12pt]{article}
\usepackage[utf8]{inputenc}
\usepackage[T2A]{fontenc}
\usepackage[russian]{babel}
\usepackage{amsmath, amssymb, amsfonts}
\usepackage{geometry}
\usepackage{multirow}
\usepackage{array}
\usepackage{graphicx, changepage}
\usepackage{pdflscape}
\usepackage{listings}
\usepackage{xcolor}
\usepackage{hyperref}
\usepackage[font=small,labelfont=bf]{caption}

\geometry{a4paper, margin=2.5cm}

\DeclareMathOperator*{\argmax}{arg\,max}
\DeclareMathOperator*{\argmin}{arg\,min}
\DeclareMathOperator*{\Argmax}{Arg\,max}
\DeclareMathOperator*{\Argmin}{Arg\,min}

\author{Лазар В. И.}
\title{Применение СДУ в задачах оценки свойств лекарственных средств}


\begin{document}
\maketitle

\newpage

\tableofcontents

\newpage

\section{Уравнение}


\[
	\begin{cases}
		dC_t = \frac{D F}{V_d} dt + (k_a - k_{el}) C_t dt + (k_a - k_{el}) \sigma dW_t, \; t \le \tau \\
		dC_t = - k_{el} C_t dt - k_{el} \sigma dW_t, \; t > \tau                                      \\
		C_0 \stackrel{\normalfont \tiny \mbox{п. н.}}{=} 0                                            \\
	\end{cases}
\]

где $W_t$ - стандартный винеровский процесс

\section{Разбиение задачи}

Разобьём данное СДУ на две последовательные задачи:

\begin{itemize}
	\item Фаза абсорбции:
	      \[
		      \begin{cases}
			      dC_t^a = \frac{DF}{V_d} dt + (k_a - k_{el}) C_t^a dt + (k_a - k_{el}) \sigma dW_t \\
			      C_0^a \stackrel{\normalfont \tiny \mbox{п. н.}}{=} 0                              \\
		      \end{cases}
	      \]

	\item Фаза элиминации:
	      \[
		      \begin{cases}
			      dC_t^{el} = - k_{el} C_t^{el} dt - k_{el} \sigma dW_{t}                \\
			      C_{\tau}^{el} \stackrel{\normalfont \tiny \mbox{п. н.}}{=}  C_{\tau}^a \\
		      \end{cases}
	      \]
\end{itemize}

\section{Решение}

\subsection{Фаза абсорбции}

Сделаем следующие замены:

\[
	\begin{cases}
		C_t^a = X_t        \\
		\frac{DF}{V_d} = A \\
		k_a - k_{el} = B   \\
		\sigma (k_a - k_{el}) = C
	\end{cases}
\]

Тогда уравнение примет следующий вид:

\[
	\begin{cases}
		dX_t = A dt + B X_t dt + C dW_t                    \\
		X_0 \stackrel{\normalfont \tiny \mbox{п. н.}}{=} 0 \\
	\end{cases}
\]

Взяв в качестве интегрирующего множителя $M_t = e^{-Bt}$ получим следующее СДУ:

\[
	d(e^{-Bt}X_t) = Ae^{-Bt}dt + Ce^{-Bt}dW_t
\]

Тогда

\[
	e^{-Bt} X_t - X_0 = A \int_0^t e^{-Bs} ds + C \int_0^t e^{-Bs} dW_s
\]

Поскольку $X_0 \stackrel{\normalfont \tiny \mbox{п. н.}}{=} 0$,

\[
	X_t = \frac{A}{B}(e^{Bt} - 1) + C \int_0^t e^{B(t-s)}dW_t
\]

Пусть $I_t = \int_0^t e^{B(t-s)}dW_t$, тогда

\[
	\mathbb{E}I_t = 0, \; \; \;
	\mathbb{D}I_t = \int_0^t e^{2B(t-s)}ds = \frac{C^2}{2B} (e^{2Bt} - 1) = D(t), \; \; \;
	I_t \sim \mathcal{N}(0, D(t))
\]

Итого:

\[
	X_t \sim \mathcal{N} (\frac{A}{B}(e^{Bt} - 1), \frac{C^2}{2B} (e^{2Bt} - 1))
\]

Сделав обратную замену, получаем выражение для $C_t^a$:

\[
	C_t^a \sim \mathcal{N} (\frac{DF}{V_d (k_a - k_{el})}(e^{(k_a - k_{el}) t} - 1), \frac{\sigma^2 (k_a - k_{el})}{2} (e^{2(k_a - k_{el})t} - 1))
\]

\subsection{Фаза элиминации}

Сделаем аналогичные замены:

\[
	\begin{cases}
		C_t^{el} = X_t     \\
		-k_{el} = A        \\
		-k_{el} \sigma = B \\
		C_{\tau}^a = x_{\tau}
	\end{cases}
\]

Получим следующее уравнение:

\[
	\begin{cases}
		dX_t = A X_t dt + B dW_t \\
		X_{\tau} = x_{\tau}
	\end{cases}
\]

Применив метод интегрирующего множителя ($M_t = e^{-At}$) и используя аналогичные рассуждения при решении получаем следующее выражение для $X_t$:

\[
	X_t = e^{A(t - \tau)} x_{\tau} + B \int_{\tau}^{t} e^{A(t - s)} dW_s
\]

Аналогично получаем для $X_t$:

\[
	X_t \sim \mathcal{N}(e^{A(t - \tau)} x_{\tau} , \frac{B^2}{2A}(e^{2A(t - \tau)} - 1))
\]

Совершаем обратную замену:

\[
	C_t^{el} \sim \mathcal{N}(e^{-k_{el}(t - \tau)} C_{\tau}^a , \frac{\sigma^2 k_{el}}{2}(1 - e^{-2k_{el}(t - \tau)})
\]

\section{Общие выводы}

В качестве решения получаем гауссовский процесс $C_t$ со следующими свойствами:

\[
	\begin{cases}
		C_t = C_t^a,    & t \le \tau \\
		C_t = C_t^{el}, & t > \tau
	\end{cases}
\]

\end{document}
